\documentclass[10pt,a4paper]{amsart}
\usepackage[margin=19mm]{geometry}
\usepackage{amsmath,amssymb,mathtools}
\usepackage[scr=rsfs,cal=euler]{mathalfa}
\usepackage{xfrac}
\usepackage[unicode,hidelinks,bookmarks=false]{hyperref}
\newcommand{\ie}{i.e.\ }
\newcommand{\cf}{cf.\ }
\newcommand{\resp}{respectively\ }
\newcommand{\Subsection}{\subsection}
\providecommand{\nobreakdash}{\nobreak\hbox{-}}
\setlength{\emergencystretch}{2em}
\multlinegap0pt
\usepackage{amssymb}
\usepackage{dsfont}
\usepackage[shortlabels]{enumitem}
\newtheorem{lem}{Lemma}
\title{Stein Variational Gradient Descent dynamics for highly concentrated kernels}
\author{Jos\'e A. Carrillo, Jakub Skrzeczkowski, Jethro Warnett}
\date{Source: arXiv:2605.03627v1 (CC BY 4.0)}
\begin{document}
\maketitle

\noindent\emph{Source and licence.} Stein Variational Gradient Descent dynamics for highly concentrated kernels. Version arXiv:2605.03627v1 (CC BY 4.0), distributed by arXiv under the Creative Commons Attribution 4.0 International licence, http://creativecommons.org/licenses/by/4.0/, as recorded in the arXiv licence metadata for this exact version and shown on the abstract page https://arxiv.org/abs/2605.03627v1. Full licence notice: \texttt{licenses/arxiv-2605.03627-CC-BY-4.0.txt}.\par\smallskip

\noindent\textit{Editorial scope.} Counted unit: the article's lem lem:optimal\_slsi\_val (source0.tex lines 1736--1746). The article's complete original proof of it is delivered with it (source0.tex lines 1747--1776, one \texttt{proof} environment of 30 lines). No mathematics is written, repaired or re-derived: the statement and the proof are the article's own lines, in the article's own notation, and no retained line is substituted or re-flowed.  No further source-local context is needed: every cross-reference of every retained line resolves inside the two delivered spans.  Nothing was omitted from any retained span, so the package carries no omission record.  No retained line contains a citation, so the wrapper carries no bibliography and no bibliography entry.  No retained line contains a figure environment, an image-inclusion command or drawing code, so no image is delivered and the package declares no asset dependency.  Editorial formatting only, in addition: an AMS \texttt{amsart} preamble that restates the article's own declarations and macros used by the retained lines, namely the environments \texttt{lem} on separate counters, together with the standard packages those lines need.  The retained environments print in this document as Lemma 1 in order of appearance, whereas the article prints the same items with the numbering of its own full document.  The complete native source of the article as distributed in the arXiv e-print for this exact version is bundled unchanged as \texttt{sources/199801/source0.tex}.
\begin{lem}
    \label{lem:optimal_slsi_val}
    Let $f\in C^0(\mathbb{R}^{d})$ and $D\geq 1$ be such that $f(0)=1$ and
    \begin{align*}
        f(x)\geq \frac{1}{D}\frac{1}{1+|x|^2}.
    \end{align*}
    For all $\sigma\in (0,1)$ define $f_\sigma(x):=f(\sigma x)$. Then for all $\varepsilon\in (0,1)$, there exists $\sigma_\ast\in (0,1)$ such that for all $\sigma\in (0,\sigma_\ast)$ we have
    \begin{align*}
        f_\sigma(x)\geq \frac{1-\varepsilon}{1+|x|^2}.
    \end{align*}
\end{lem}

\begin{proof}
    Let $c := D(1-\varepsilon)$. Since $f$ is continuous at $0$ and $f(0)=1$, there exists $\delta > 0$ such that
    \begin{align*}
        \inf_{|y|< \delta}f(y) \geq 1 - \varepsilon.
    \end{align*}
    Now, choose $\sigma_\ast \in (0,1)$ sufficiently small such that for all $\sigma \in (0,\sigma_\ast)$
    \begin{align}
        \label{eq:sigma_choice}
        c \sigma^2 < 1
        \qquad \text{and}\qquad1 + \frac{\delta^2}{\sigma^2} \geq c(1+\delta^2).
    \end{align}
    Next we split into two cases. The first case is when $|x| < \delta/\sigma$. Then $|\sigma x| < \delta$, so
    \begin{align*}
        f_\sigma(x)=f(\sigma x) \geq 1 - \varepsilon
        \geq \frac{1-\varepsilon}{1+|x|^2}.
    \end{align*}
    It remains to consider the other case $|x| \geq \delta/\sigma$. Using the assumed bound on $f$, we have
    \begin{align*}
        f(\sigma x) \geq \frac{1}{D(1+\sigma^2|x|^2)}.
    \end{align*}
    Thus, it suffices to show
    \begin{align*}
        \frac{1}{D(1+\sigma^2|x|^2)} \geq \frac{1-\varepsilon}{1+|x|^2} \iff (1-c) + (1-c\,\sigma^2)\,|x|^2 \geq 0.
    \end{align*}
    Note that $(1-c\,\sigma^2) \geq 0$ by \eqref{eq:sigma_choice} so we can estimate using $|x|^2 \geq \delta^2/\sigma^2$
    $$
    (1-c) + (1-c\,\sigma^2)\,|x|^2 \geq (1-c) + (1-c\,\sigma^2)\,\frac{\delta^2}{\sigma^2}  = 1 + \frac{\delta^2}{\sigma^2} - c\,(1+\delta^2)\geq 0
    $$
    by the second item in \eqref{eq:sigma_choice}. The proof is concluded.
\end{proof}

\end{document}
